\documentclass[a4paper,12pt]{article}
\usepackage[utf8]{inputenc}
\usepackage[portuguese]{babel}
\usepackage{graphicx}
\usepackage{amsmath}
\usepackage{enumitem}
\usepackage{natbib}
\usepackage{indentfirst}
\usepackage{hyperref}
\usepackage{fancyhdr}
\usepackage{listings}
\usepackage{xcolor}
\usepackage{acronym}
\usepackage{dirtree}

\lstset{
    language=Java,
    basicstyle=\ttfamily\small,
    keywordstyle=\color{blue},
    stringstyle=\color{red},
    commentstyle=\color{green!50!black},
    numberstyle=\tiny,
    numbers=left,
    stepnumber=1,
    numbersep=5pt,
    backgroundcolor=\color{gray!10},
    showspaces=false,
    showstringspaces=false,
    showtabs=false,
    tabsize=2,
    captionpos=b,
    breaklines=true,
    breakatwhitespace=true,
    breakautoindent=true,
    escapeinside={\%*}{*)},
    linewidth=\textwidth,
    basewidth=0.5em,
}

\pagestyle{fancy}
\fancyhf{}
\fancyfoot[R]{\thepage}
\renewcommand{\headrulewidth}{0pt}

\begin{document}

% Logo da universidade
\begin{figure}
\centering
\includegraphics{Anexos/UBI_LOGO.jpg}
\end{figure}

% Título e autores
\title{Relatório do Trabalho Prático 2 - Farmácia GUI}
\author{Afonso Esteves - 55047\\Felipe Dourado - 53394\\Tiago Pereira - 55019}
\date{\today}

\maketitle
\newpage

% Sumário
\tableofcontents

\newpage

\section*{Resumo}
\hspace{10pt}
Este relatório documenta a evolução de uma aplicação de gestão de farmácia, migrando de uma interface de consola para uma interface gráfica de utilizador (GUI) desenvolvida com Java Swing. O sistema foi reestruturado em pacotes (modelo, GUI, persistência) para melhor organização e manutenibilidade, seguindo princípios de arquitetura em camadas. Além da nova interface, foi implementada a persistência de dados através de serialização de objetos, permitindo que o estado da aplicação (clientes, produtos e encomendas) seja guardado e recuperado entre sessões. Esta segunda versão representa um avanço significativo em termos de usabilidade e robustez.
\addcontentsline{toc}{section}{Resumo}
\section*{Palavras-chave e \textit{Keywords}}
Java; Swing; GUI; Programação Orientada a Objetos; Persistência de Dados; Serialização; JUnit; Javadoc.

\newpage

\section{Introdução}
\subsection{Âmbito, enquadramento e motivação}
\hspace{10pt}
Este documento detalha o desenvolvimento do Trabalho Prático 2, uma evolução do sistema de gestão de farmácia introduzido no TP1. O foco principal desta nova fase foi a substituição da interface de linha de comandos por uma interface gráfica (GUI) mais intuitiva e funcional, utilizando a biblioteca Java Swing. Adicionalmente, foi introduzida uma camada de persistência para garantir que os dados da aplicação não são perdidos ao fechar o programa, e foram adicionados testes unitários e documentação formal ao projeto.

A motivação foi aprofundar os conhecimentos em Programação Orientada a Objetos, explorando conceitos mais avançados como o desenvolvimento de interfaces gráficas, programação orientada a eventos, serialização de objetos, testes com JUnit, documentação com Javadoc e a importância de uma arquitetura de software bem definida.

\section{Contextualização do Projeto}

\subsection{Identificação dos Desenvolvedores}
Os autores do presente projeto são estudantes do curso de Informática Web, Móvel e na Nuvem da \ac{UBI}, matriculados na disciplina de Programação Orientada a Objetos.

\subsubsection{Desenvolvedores Envolvidos}
\begin{itemize}
    \item \textbf{Afonso Esteves (55047)}
    \item \textbf{Felipe Dourado (53394)}
    \item \textbf{Tiago Pereira (55019)}
\end{itemize}     

\subsection{Contextualização Académica}
O projeto foi desenvolvido no contexto da disciplina de Programação Orientada a Objetos, lecionada pelo Professor Paul Andrew Crocker, com o objetivo de aplicar e consolidar conhecimentos em:
\begin{itemize}
    \item Desenvolvimento de Interfaces Gráficas (GUI) com Java Swing.
    \item Programação Orientada a Eventos (ActionListeners, MouseListeners).
    \item Persistência de dados através de serialização de objetos Java.
    \item Arquitetura de software em camadas (Modelo, Visão, Persistência).
    \item Testes unitários com JUnit para garantir a qualidade do código.
    \item Geração de documentação técnica com Javadoc.
    \item Refatoração e melhoria de código existente.
\end{itemize}

\section{Especificações Técnicas}

\subsection{Ambiente de Desenvolvimento}
\begin{itemize}
    \item \textbf{Linguagem:} Java (JDK 8 ou superior)
    \item \textbf{Bibliotecas Gráficas:} Java Swing
    \item \textbf{Testes:} JUnit 5
    \item \textbf{Documentação:} Javadoc
    \item \textbf{Ambiente de Execução:} Java Development Kit (JDK)
    \item \textbf{Ambiente de Desenvolvimento:} IntelliJ IDEA
\end{itemize}

\section{Utilização}
\subsection{Requisitos Prévios}
\begin{itemize}
    \item \textbf{Java Development Kit (JDK) 8 ou superior:}
    Essencial para compilar e executar a aplicação. Pode ser descarregado do site oficial da Oracle.
\end{itemize}

\subsection{Compilação e Execução}
\begin{enumerate}
    \item \textbf{Navegar até ao diretório raiz do projeto.}
    \item \textbf{Compilar os ficheiros Java:}
    \begin{verbatim}
javac -d TP-02/out/production/TP-02 -cp TP-02/lib/junit-platform-console-standalone-1.10.2.jar $(find TP-02/src -name "*.java")
    \end{verbatim}
    \item \textbf{Executar a aplicação GUI:}
    \begin{verbatim}
java -cp TP-02/out/production/TP-02 Main
    \end{verbatim}
\end{enumerate}

\subsection{Execução de Testes Unitários}
Para verificar a integridade da lógica de negócio, podem ser executados os testes JUnit 5:
\begin{enumerate}
    \item \textbf{Navegar até ao diretório raiz do projeto.}
    \item \textbf{Executar os testes (após compilação):}
    \begin{verbatim}
java -jar TP-02/lib/junit-platform-console-standalone-1.10.2.jar -cp TP-02/out/production/TP-02 --scan-classpath
    \end{verbatim}
\end{enumerate}

\subsection{Geração de Documentação Javadoc}
Para gerar a documentação HTML do código fonte:
\begin{enumerate}
    \item \textbf{Navegar até ao diretório raiz do projeto.}
    \item \textbf{Gerar a documentação:}
    \begin{verbatim}
javadoc -d TP-02/doc TP-02/src/**/*.java TP-02/src/*.java
    \end{verbatim}
    A documentação ficará disponível no diretório `TP-02/doc/`.
\end{enumerate}

\section{Estrutura da Solução Construída}
O sistema foi reestruturado para uma arquitetura em camadas, separando as responsabilidades em pacotes distintos.

\subsection{Componentes do Sistema}
\begin{itemize}
    \item \textbf{Pacote `modelo`:} Contém as classes que representam o domínio do problema (\texttt{Cliente}, \texttt{Produto}, \texttt{Farmacia}, etc.).
    \item \textbf{Pacote `gui`:} Contém as classes da interface gráfica (\texttt{FarmaciaGUI}, \texttt{PainelGestaoClientes}, etc.).
    \item \textbf{Pacote `persistencia`:} Contém a classe \texttt{GestorDados}, responsável por guardar e carregar o estado da aplicação.
    \item \textbf{Pacote `test`:} Contém as classes de teste unitário, como \texttt{FarmaciaTest.java}.
    \item \textbf{Classe `Main`:} Ponto de entrada que inicializa a farmácia, carrega os dados e lança a GUI.
\end{itemize}

\subsection{Testes e Documentação}
O projeto inclui artefactos importantes para garantir a qualidade e a manutenibilidade do software:
\begin{itemize}
    \item \textbf{Testes Unitários:} O diretório `TP-02/src/test/` contém a classe `FarmaciaTest.java`, que utiliza o framework JUnit 5 para validar a lógica de negócio da classe `Farmacia`, como adicionar clientes, gerir stock e processar encomendas.
    \item \textbf{Documentação Javadoc:} O diretório `TP-02/doc/` é o destino da documentação gerada pela ferramenta Javadoc. Esta documentação descreve as classes, métodos e a API do projeto, facilitando a sua compreensão e reutilização.
\end{itemize}

\section{Funcionalidades Implementadas}
As funcionalidades do sistema foram expandidas e agora são acedidas através da interface gráfica:
\begin{itemize}
    \item \textbf{Interface Gráfica com Separadores:} A janela principal usa um \texttt{JTabbedPane} para separar a gestão de Clientes, Produtos e Encomendas.
    \item \textbf{Gestão CRUD Completa:} Todas as entidades podem ser criadas, lidas, atualizadas e removidas (CRUD) através de formulários e tabelas interativas.
    \item \textbf{Persistência Automática:} Os dados da aplicação são automaticamente carregados ao iniciar e guardados ao fechar a janela.
    \item \textbf{Testes Unitários:} Cobertura de testes para a lógica de negócio principal, garantindo robustez.
\end{itemize}

\section{Implementação de Funcionalidades}

\subsection{Classe FarmaciaGUI: A Janela Principal}
A classe `FarmaciaGUI` configura a janela principal, o `JTabbedPane` e o comportamento de fecho para guardar os dados.

\begin{minipage}{\linewidth}
\begin{lstlisting}[caption={Configuração da janela principal em FarmaciaGUI.java}]
public class FarmaciaGUI extends JFrame {
    // ...
    public FarmaciaGUI(Farmacia farmacia) {
        // ...
        setTitle("Gestão de Farmácia");
        setSize(800, 600);
        setDefaultCloseOperation(JFrame.DO_NOTHING_ON_CLOSE);
        setLocationRelativeTo(null);

        addWindowListener(new WindowAdapter() {
            @Override
            public void windowClosing(WindowEvent e) {
                // ... Lógica para guardar dados ...
            }
        });

        JTabbedPane tabbedPane = new JTabbedPane();
        // ... criação e adição dos painéis ...
        add(tabbedPane);
    }
}
\end{lstlisting}
\end{minipage}

\subsection{Classe GestorDados: Persistência de Dados}
A classe `GestorDados` usa serialização para guardar e carregar as listas de objetos que compõem o estado da aplicação.

\begin{minipage}{\linewidth}
\begin{lstlisting}[caption={Método para guardar dados em GestorDados.java}]
public static void guardarDados(List<Cliente> clientes, List<Produto> produtos, List<Encomenda> encomendas) {
    try (ObjectOutputStream oos = new ObjectOutputStream(new FileOutputStream(NOME_FICHEIRO))) {
        oos.writeObject(clientes);
        oos.writeObject(produtos);
        oos.writeObject(encomendas);
    } catch (IOException e) {
        // ...
    }
}
\end{lstlisting}
\end{minipage}

\subsection{Classe FarmaciaTest: Testes Unitários}
Os testes com JUnit 5 validam o comportamento esperado da lógica de negócio, como a adição de um cliente.

\begin{minipage}{\linewidth}
\begin{lstlisting}[caption={Exemplo de um teste unitário em FarmaciaTest.java}]
class FarmaciaTest {
    private Farmacia farmacia;

    @BeforeEach
    void setUp() {
        farmacia = new Farmacia();
    }

    @Test
    void testAdicionarCliente() {
        Cliente cliente = new Cliente("Nome Teste", "123456789", new Morada("Rua Teste", "Cidade Teste", "6200-000"));
        farmacia.adicionarCliente(cliente);
        assertEquals(1, farmacia.getClientes().size());
        assertEquals("Nome Teste", farmacia.acharClienteNif("123456789").get().getNome());
    }
}
\end{lstlisting}
\end{minipage}

\section{Considerações de Implementação}

\subsection{Princípios de Desenvolvimento}
\begin{itemize}
    \item \textbf{Arquitetura MVC (Model-View-Controller):} A estrutura em pacotes `modelo` (Model), `gui` (View) e a lógica de eventos nos painéis (Controller) aproxima o projeto do padrão MVC.
    \item \textbf{Programação Orientada a Eventos:} A GUI é controlada por eventos gerados pelo utilizador, tratados por `Listeners`.
    \item \textbf{Test-Driven Development (TDD):} A inclusão de testes unitários com JUnit incentiva um ciclo de desenvolvimento onde os testes guiam a implementação da lógica de negócio.
\end{itemize}

\section{Conclusão}
O Trabalho Prático 2 representa um marco importante na evolução do sistema de farmácia. A migração para uma interface gráfica com Swing melhorou a usabilidade, enquanto a implementação da persistência de dados com serialização tornou a aplicação mais robusta. A adição de testes unitários com JUnit e a geração de documentação com Javadoc elevaram a qualidade e a profissionalismo do projeto, solidificando os conhecimentos em engenharia de software e Programação Orientada a Objetos.

\clearpage

\section{Referências Bibliográficas}
\begin{itemize}
    \item Schildt, Herbert. \textit{Java: A Beginner's Guide}. 6th ed., McGraw-Hill Education, 2014.
    \item Olsson, Mikael. \textit{Java Quick Syntax Reference}. Apress, 2014.
    \item Roberts, Eric S. \textit{The Art and Science of Java}. Addison-Wesley, 2008.
\end{itemize}

\section*{Acrónimos}
\begin{acronym}[UBI]
    \acro{UBI}{Universidade da Beira Interior}
\end{acronym}

\end{document}

